% layout: split
% font: footnotesize
% TAB-01: generated by harness/analysis/paper_outputs.py; do not edit (every cell's source is in TAB-01.cells.json)
\begin{tabularx}{\linewidth}{llrr>{\hsize=1.583\hsize\raggedright\hyphenpenalty=10000\exhyphenpenalty=10000\arraybackslash}X>{\hsize=0.417\hsize\raggedleft\hyphenpenalty=10000\exhyphenpenalty=10000\arraybackslash}X}
\toprule
Test & Study & Puzzles & Steps & What is compared & Margin (points) \\
\midrule
Test~1 & Study~1 & 80 & 6 & Haiku-medium vote vs Opus-medium, 6\ensuremath{\times}1 & \ensuremath{-}10 \\
Test~2 & Study~1 & 80 & 6 & Haiku-medium vote vs Opus-medium, 1\ensuremath{\times}6 & \ensuremath{-}10 \\
Test~3 & Study~1 & 80 & 6 & Haiku-medium vote vs Sonnet-medium, 1\ensuremath{\times}6 & \ensuremath{-}10 \\
Test~4 & Study~1 & 80 & 6 & Step voting vs Opus-medium, 1\ensuremath{\times}6 & \ensuremath{-}10 \\
Test~5 & Study~1 & 80 & 6 & Opus-low vs Sonnet-medium, 1\ensuremath{\times}6 & \ensuremath{-}10 \\
\addlinespace
B1 & Study~1b & 80 & 16 & One Haiku-max sample vs one Opus-medium answer & \ensuremath{-}5 \\
B2 & Study~1b & 80 & 16 & Haiku-max vote vs one Opus-medium answer & \ensuremath{-}5 \\
B3 & Study~1b & 40 & 6 & One Haiku-max sample vs one Opus-medium answer & \ensuremath{-}5 \\
\bottomrule
\end{tabularx}
